\documentclass{amsart}
% For dealing with references we use the comment environment
\usepackage{verbatim}
\newenvironment{reference}{\comment}{\endcomment}
%\newenvironment{reference}{}{}
\newenvironment{slogan}{\comment}{\endcomment}
\newenvironment{history}{\comment}{\endcomment}

% For commutative diagrams we use Xy-pic
\usepackage[all]{xy}

% We use 2cell for 2-commutative diagrams.
\xyoption{2cell}
\UseAllTwocells

% We use multicol for the list of chapters between chapters
\usepackage{multicol}

% This is generally recommended for better output
\usepackage{lmodern}
\usepackage[T1]{fontenc}

% For cross-file-references

% Package for hypertext links:
\usepackage{hyperref}

% For any local file, say "hello.tex" you want to link to please

% Theorem environments.
%
\theoremstyle{plain}
\newtheorem{theorem}[subsection]{Theorem}
\newtheorem{proposition}[subsection]{Proposition}
\newtheorem{lemma}[subsection]{Lemma}

\theoremstyle{definition}
\newtheorem{definition}[subsection]{Definition}
\newtheorem{example}[subsection]{Example}
\newtheorem{exercise}[subsection]{Exercise}
\newtheorem{situation}[subsection]{Situation}

\theoremstyle{remark}
\newtheorem{remark}[subsection]{Remark}
\newtheorem{remarks}[subsection]{Remarks}

\numberwithin{equation}{subsection}

% Macros
%
\def\lim{\mathop{\mathrm{lim}}\nolimits}
\def\colim{\mathop{\mathrm{colim}}\nolimits}
\def\Spec{\mathop{\mathrm{Spec}}}
\def\Hom{\mathop{\mathrm{Hom}}\nolimits}
\def\Ext{\mathop{\mathrm{Ext}}\nolimits}
\def\SheafHom{\mathop{\mathcal{H}\!\mathit{om}}\nolimits}
\def\SheafExt{\mathop{\mathcal{E}\!\mathit{xt}}\nolimits}
\def\Sch{\mathit{Sch}}
\def\Mor{\mathop{\mathrm{Mor}}\nolimits}
\def\Ob{\mathop{\mathrm{Ob}}\nolimits}
\def\Sh{\mathop{\mathit{Sh}}\nolimits}
\def\NL{\mathop{N\!L}\nolimits}
\def\CH{\mathop{\mathrm{CH}}\nolimits}
\def\proetale{{pro\text{-}\acute{e}tale}}
\def\etale{{\acute{e}tale}}
\def\QCoh{\mathit{QCoh}}
\def\Ker{\mathop{\mathrm{Ker}}}
\def\Im{\mathop{\mathrm{Im}}}
\def\Coker{\mathop{\mathrm{Coker}}}
\def\Coim{\mathop{\mathrm{Coim}}}

% Boxtimes
%
\DeclareMathSymbol{\boxtimes}{\mathbin}{AMSa}{"02}

%
% Macros for moduli stacks/spaces
%
\def\QCohstack{\mathcal{QC}\!\mathit{oh}}
\def\Cohstack{\mathcal{C}\!\mathit{oh}}
\def\Spacesstack{\mathcal{S}\!\mathit{paces}}
\def\Quotfunctor{\mathrm{Quot}}
\def\Hilbfunctor{\mathrm{Hilb}}
\def\Curvesstack{\mathcal{C}\!\mathit{urves}}
\def\Polarizedstack{\mathcal{P}\!\mathit{olarized}}
\def\Complexesstack{\mathcal{C}\!\mathit{omplexes}}
% \Pic is the operator that assigns to X its picard group, usage \Pic(X)
% \Picardstack_{X/B} denotes the Picard stack of X over B
% \Picardfunctor_{X/B} denotes the Picard functor of X over B
\def\Pic{\mathop{\mathrm{Pic}}\nolimits}
\def\Picardstack{\mathcal{P}\!\mathit{ic}}
\def\Picardfunctor{\mathrm{Pic}}
\def\Deformationcategory{\mathcal{D}\!\mathit{ef}}

\setlength{\emergencystretch}{3em}
\begin{document}
\title[Stacks Project, Tag 0G5E]{266 --- Nonvanishing of top coherent cohomology forces a variety to be proper}
\maketitle
\noindent Copyright (C) 2005--2025 Johan de Jong. Stacks Project authors. Source: \href{https://stacks.math.columbia.edu/tag/0G5E}{Tag 0G5E}.

\noindent Editorial difficulty note (not author text). Difficulty H2: The proof compactifies the variety and converts top cohomology into compactly supported cohomology of its dual. It proves the relevant lowest dual cohomology sheaf is torsion-free and consequently has no section supported away from a nonempty boundary. This is the key duality-and-support vanishing argument establishing properness..

\noindent Editorial provenance note (not author text). History: excerpt from revision \texttt{a0444\allowbreak 6e57e\allowbreak c1fbc\allowbreak 252a8\allowbreak 71afc\allowbreak ec775\allowbreak 2fb28\allowbreak 07b14}, duality.tex, lines 9111--9161. Reference presentation was adapted; original-author context and core support blocks were retained or added. The mathematical wording is unchanged. Licensed under GNU FDL 1.2 or later; no invariant sections or cover texts. See ../licenses/COPYING. Context blocks reproduce author definitions and supporting results; they are not additional counted problems.

\section{Duality for compactly supported cohomology}
\label{section-duality-compactly-supported}

\noindent
Let $k$ be a field. Let $U$ be a separated scheme of finite type over $k$.
Let $K$ be an object of $D^b_{\textit{Coh}}(\mathcal{O}_U)$. Let us define
the compactly supported cohomology $H^i_c(U, K)$ of $K$ as follows.
Choose an open immersion $j : U \to X$ into a scheme proper over $k$
and a Deligne system $(K_n)$ for $j : U \to X$ whose restriction
to $U$ is constant with value $K$. Then we set
$$
H^i_c(U, K) = \lim H^i(X, K_n)
$$
We view this as a topological $k$-vector space using the limit topology
(see More on Algebra, Section \href{https://stacks.math.columbia.edu/tag/07E7}{section topological ring (Tag 07E7)}).
There are several points to make here.

\medskip\noindent
First, this definition is independent of the choice of $X$ and $(K_n)$.
Namely, if $p : U \to \Spec(k)$ denotes the structure morphism, then
we already know that $Rp_!K = (R\Gamma(X, K_n))$
is well defined up to pro-isomorphism in $D(k)$ hence so is the limit
defining $H^i_c(U, K)$.

\medskip\noindent
Second, it may seem more natural to use the expression
$$
H^i(R\lim R\Gamma(X, K_n)) = R\Gamma(X, R\lim K_n)
$$
but this would give the same answer: since the $k$-vector spaces
$H^j(X, K_n)$ are finite dimensional, these inverse systems satisfy
Mittag-Leffler and hence $R^1\lim$ terms of Cohomology, Lemma
\href{https://stacks.math.columbia.edu/tag/0D60}{lemma RGamma commutes with Rlim (Tag 0D60)} vanish.

\medskip\noindent
If $U' \subset U$ is an open subscheme, then there is a canonical map
$$
H^i_c(U', K|_{U'}) \longrightarrow H^i_c(U, K)
$$
functorial for $K$ in $D^b_{\textit{Coh}}(\mathcal{O}_U)$.
See for example Remark \href{https://stacks.math.columbia.edu/tag/0G56}{remark covariance open lower shriek (Tag 0G56)}.
In fact, using Remark \href{https://stacks.math.columbia.edu/tag/0G57}{remark covariance etale lower shriek (Tag 0G57)}
we see that more generally such a map exists for an \'etale morphism
$U' \to U$ of separated schemes of finite type over $k$.

\medskip\noindent
If $V$ is a $k$-vector space then we put a topology on $\Hom_k(V, k)$
as follows: write $V = \bigcup V_i$ as the filtered union of its finite
dimensional $k$-subvector spaces and use the limit topology on
$\Hom_k(V, k) = \lim \Hom_k(V_i, k)$. If $\dim_k V < \infty$ then
the topology on $\Hom_k(V, k)$ is discrete. More generally, if
$V = \colim_n V_n$ is written as a directed colimit of finite dimensional
vector spaces, then $\Hom_k(V, k) = \lim \Hom_k(V_n, k)$ as topological
vector spaces.



\begin{lemma}
\label{lemma-lichtenbaum}
Let $U$ be a variety. Let $\mathcal{F}$ be a coherent $\mathcal{O}_U$-module.
If $H^d(U, \mathcal{F})$ is nonzero, then $\dim(U) \geq d$ and if
equality holds, then $U$ is proper.
\end{lemma}

\begin{proof}
By the Grothendieck's vanishing result in
Cohomology, Proposition \href{https://stacks.math.columbia.edu/tag/02UZ}{proposition vanishing Noetherian (Tag 02UZ)}
we conclude that $\dim(U) \geq d$.
Assume $\dim(U) = d$. Choose a compactification $U \to X$
such that $U$ is dense in $X$. (This is possible by
More on Flatness, Theorem \href{https://stacks.math.columbia.edu/tag/0F41}{theorem nagata (Tag 0F41)} and
Lemma \href{https://stacks.math.columbia.edu/tag/0A9Z}{lemma compactifyable (Tag 0A9Z)}.) After replacing $X$ by its
reduction we find that $X$ is a proper variety of dimension $d$
and we see that $U$ is proper if and only if $U = X$.
Set $Z = X \setminus U$. We will show that $H^d(U, \mathcal{F})$
is zero if $Z$ is nonempty.

\medskip\noindent
Choose a coherent $\mathcal{O}_X$-module
$\mathcal{G}$ whose restriction to $U$ is $\mathcal{F}$, see
Properties, Lemma \href{https://stacks.math.columbia.edu/tag/0G41}{lemma lift finite presentation (Tag 0G41)}.
Let $\omega_X^\bullet$ denote the dualizing complex of $X$ as in
Section \href{https://stacks.math.columbia.edu/tag/0FVU}{section duality proper over field (Tag 0FVU)}.
Set $\omega_U^\bullet = \omega_X^\bullet|_U$.
Then $H^d(U, \mathcal{F})$ is dual to
$$
H^{-d}_c(U, R\SheafHom_{\mathcal{O}_U}(\mathcal{F}, \omega_U^\bullet))
$$
by Lemma \href{https://stacks.math.columbia.edu/tag/0G5A}{lemma duality compact support (Tag 0G5A)}. By
Lemma \href{https://stacks.math.columbia.edu/tag/0FVV}{lemma duality proper over field (Tag 0FVV)} we see that
the cohomology sheaves of $\omega_X^\bullet$ vanish in degrees $< -d$
and $H^{-d}(\omega_X^\bullet) = \omega_X$ is a coherent
$\mathcal{O}_X$-module which is $(S_2)$ and whose support is $X$.
In particular, $\omega_X$ is torsion free, see
Divisors, Lemma \href{https://stacks.math.columbia.edu/tag/0AXY}{lemma torsion free finite noetherian domain (Tag 0AXY)}.
Thus we see that the cohomology sheaf
$$
H^{-d}(R\SheafHom_{\mathcal{O}_X}(\mathcal{G}, \omega_X^\bullet)) =
\SheafHom(\mathcal{G}, \omega_X)
$$
is torsion free, see
Divisors, Lemma \href{https://stacks.math.columbia.edu/tag/0AXZ}{lemma hom into torsion free (Tag 0AXZ)}.
Consequently this sheaf has no nonzero sections vanishing
on any nonempty open of $X$ (those would be torsion sections).
Thus it follows from Lemma \href{https://stacks.math.columbia.edu/tag/0G5C}{lemma h0 compactly supported (Tag 0G5C)} that
$H^{-d}_c(U, R\SheafHom_{\mathcal{O}_U}(\mathcal{F}, \omega_U^\bullet))$
is zero, and hence $H^d(U, \mathcal{F})$ is zero as desired.
\end{proof}
\end{document}
